% PluralTree — GRPO alignment phase, method section.
% Standalone-compilable; the body (Sections 1–6) lifts directly into a paper.
%   pdflatex grpo_method.tex
\documentclass[11pt]{article}
\usepackage[margin=1in]{geometry}
\usepackage{amsmath,amssymb}
\usepackage{booktabs}
\usepackage[colorlinks=true,linkcolor=blue,citecolor=blue]{hyperref}

\newcommand{\pol}{\pi_\theta}
\newcommand{\polold}{\pi_{\theta_{\mathrm{old}}}}
\newcommand{\polref}{\pi_{\mathrm{ref}}}
\newcommand{\Pos}{\mathcal{P}}
\newcommand{\Ind}[1]{\mathbf{1}\!\left[#1\right]}

\title{Pluralistic Alignment via Group-Relative Policy Optimization\\
{\large with a Graph-Grounded Coverage Reward}}
\author{PluralTree}
\date{}

\begin{document}
\maketitle

\section{Setup and notation}

Let $q$ be an open-ended question on which people genuinely disagree. The
knowledge graph $\mathcal{G}$ is a rooted tree over survey data
(\textsc{root} $\to$ topic $\to$ subtopic $\to$ question $\to$ axis $\to$
opinion), where an \emph{axis} is a (question $\times$ demographic attribute)
pair and each opinion leaf carries one subgroup's probability distribution over
that question's answer options.

A retrieval module (the \emph{divergence scout}) selects an anchor
$a(q)\in\mathcal{G}$ and renders its subtree into a context block, giving the
policy input $x_q$. Retrieval is executed once, at dataset-construction time:
the graph encoder and the scout are \textbf{frozen} throughout this phase, so
what is trained is the reasoner's use of retrieved evidence, not the retrieval
itself.

\section{Graph-grounded positions}

For anchor $a(q)$ we aggregate the opinion leaves of its subtree into a set of
population-level \emph{positions}
\[
  \Pos(q)=\{(p_1,\pi_1),\dots,(p_n,\pi_n)\},
  \qquad
  \pi_k=\frac{1}{|S_k|}\sum_{g\in S_k} \Pr\nolimits_g(\text{option }k),
\]
where $S_k$ is the set of subgroups answering the underlying survey item and
$\pi_k$ is the mean subgroup probability of option $k$ (its \emph{prevalence}).
Positions below a prevalence floor are discarded as noise.

This is the crux of the design: the reward is defined against \emph{measured
human survey distributions}, not against a judge model. It therefore (i)
introduces no dependence on the held-out evaluation benchmark, and (ii) can be
computed on any question with a graph anchor, giving effectively unlimited
training signal.

\section{Reward}

Let $y$ be a generated answer, segmented into position-units
$u_1,\dots,u_m$ (enumerated items or sentences). Let $e(\cdot)$ be a
\emph{held-out} sentence embedder --- deliberately distinct from the encoder
the scout retrieves with, so the reward cannot credit merely
retrieval-shaped text --- and write
$s_{jk}=\langle e(u_j),\,e(p_k)\rangle$ for unit-normalised embeddings.

\paragraph{Mention.} Position $k$ is \emph{mentioned} if some unit matches it:
\[
  \textsc{men}_k(y)=\Ind{\max_j s_{jk}\ \ge\ \tau }.
\]

\paragraph{Depth.} Each unit is attributed to its single best-matching position,
$\alpha(j)=\arg\max_k s_{jk}$, and contributes its length only if that match
clears $\tau$:
\[
  \mathrm{dep}_k(y)=\sum_{j\,:\,\alpha(j)=k,\ s_{j\alpha(j)}\ge\tau} |u_j| ,
\]
with $|u_j|$ the word count of unit $j$. Note $\textsc{men}$ and $\mathrm{dep}$
use different assignment rules: a position may be \emph{mentioned} by a unit
whose best match lies elsewhere.

\paragraph{Coverage.} A position counts only when it is both named and
articulated,
\begin{equation}\label{eq:covered}
  \textsc{cov}_k(y)\;=\;\textsc{men}_k(y)\;\wedge\;\Ind{\mathrm{dep}_k(y)\ \ge\ d},
\end{equation}
for a depth threshold $d$. Equation~\eqref{eq:covered} is the central
correction relative to a mention-only reward. Empirically, an answer that names
every position in a single clause each scores near zero on the target metric
while a mention-only reward scores it as full recall; without the depth term the
reward is blind to the dominant failure mode.

\paragraph{Terms.} With $\tau$ the match threshold and $T$ an effective position
count,
\begin{align}
  R(y) &= \frac{1}{n}\sum_{k=1}^{n}\textsc{cov}_k(y)
        &&\text{(recall, \emph{uniform} over positions)}\\
  P(y) &= \frac{1}{m}\sum_{j=1}^{m}\Ind{\max_k s_{jk}\ge\tau}
        &&\text{(precision; penalises padding)}\\
  V(y) &= \frac{\max\!\big(0,\ \textstyle\sum_k \textsc{cov}_k(y)-T\big)}{T},
        \quad T=\exp\!\big(H(\pi)\big)
        &&\text{(over-enumeration)}
\end{align}

Recall is \emph{unweighted} even though prevalences $\pi_k$ are available. The
evaluation metric counts every viewpoint cluster equally; weighting recall by
prevalence would train the policy to serve the majority and neglect exactly the
minority positions pluralism exists to represent.

\paragraph{Composition.} The terms combine multiplicatively,
\begin{equation}\label{eq:reward}
  r(y)\;=\;R(y)\;\cdot\;P(y)^{\lambda_p}\;\cdot\;\frac{1}{1+\lambda_v V(y)}
  \;\in[0,1],
\end{equation}
rather than as subtracted penalties. Additive penalties admit a compensatory
effect --- the policy inflates one term to pay for another --- whereas factors in
$[0,1]$ cannot be traded off, and \eqref{eq:reward} needs no clamping.
We set $\lambda_v=0$ by default: the evaluation metric is monotone in the number
of clusters covered, so it cannot reward commitment, and $T$ is derived from
\emph{graph} prevalence skew, which does not imply few \emph{human} viewpoint
clusters. The anti-enumeration role is instead served by the depth term in
\eqref{eq:covered}: listing many positions shallowly leaves none above $d$.

\section{Objective}

For each prompt $x_q$ we sample a group of $G$ completions
$\{y_i\}_{i=1}^{G}\sim\polold(\cdot\mid x_q)$ and score each with
\eqref{eq:reward}. The group supplies its own baseline, so no value network is
required:
\begin{equation}\label{eq:adv}
  A_i \;=\; \frac{r(y_i)-\mu}{\sigma+\varepsilon},
  \qquad
  \mu=\tfrac{1}{G}\textstyle\sum_i r(y_i),\quad
  \sigma=\mathrm{std}\big(r(y_1),\dots,r(y_G)\big).
\end{equation}
Writing $\rho_{i,t}(\theta)=\dfrac{\pol(y_{i,t}\mid x_q,y_{i,<t})}
{\polold(y_{i,t}\mid x_q,y_{i,<t})}$, the objective is the clipped
policy-gradient surrogate with a KL anchor to the reference policy:
\begin{equation}\label{eq:grpo}
  \mathcal{J}(\theta)=
  \mathbb{E}\Bigg[\frac{1}{G}\sum_{i=1}^{G}\frac{1}{|y_i|}\sum_{t=1}^{|y_i|}
  \min\Big(\rho_{i,t}A_i,\ \mathrm{clip}(\rho_{i,t},1-\epsilon,1+\epsilon)A_i\Big)
  \Bigg]
  \;-\;\beta\,\mathbb{D}_{\mathrm{KL}}\!\big[\pol \,\|\, \polref\big].
\end{equation}

A consequence of \eqref{eq:adv} worth stating: if every rollout in a group
receives the same reward, $\sigma\to 0$ and the group contributes no gradient.
This is benign for the per-response reward \eqref{eq:reward}, whose value
depends on response content, but it becomes a real hazard for any
\emph{group-dependent} reward (e.g.\ crediting each rollout by its marginal
contribution to the group's union coverage), where perfect complementarity and
total collapse both yield symmetric rewards and hence a vanishing advantage.

\section{Policy}

$\pol$ is a LoRA adapter over a trainable bf16 instruction-tuned base
($\sim$7B); the evaluation-time reasoner is a larger quantised model and is not
trainable. Consequently any comparison must include a \emph{same-base}
no-training control: a policy may improve over its own base and still sit below
the larger model's untrained baseline.

\section{Leakage protocol}

The evaluation benchmark's questions and its human ratings are held out
absolutely. Training prompts are drawn from the graph's own survey questions,
and any question whose normalised text matches a benchmark item is dropped at
dataset-construction time. The reward \eqref{eq:reward} never consults human
ratings or a judge model, so no evaluation signal can enter training through the
reward path.

\section{Hyperparameters}

\begin{table}[h]
\centering
\begin{tabular}{@{}lll@{}}
\toprule
Symbol & Meaning & Value \\
\midrule
$G$ & rollouts per prompt (group size) & 8 \\
$\beta$ & KL coefficient to $\polref$ & 0.04 \\
$\epsilon$ & PPO clip range & 0.2 \\
$\tau$ & match threshold (cosine) & 0.50 \\
$d$ & depth threshold (words/position) & 60$^{\dagger}$ \\
$\lambda_p$ & precision exponent & 0.5 \\
$\lambda_v$ & over-enumeration scale & 0 \\
-- & sampling temperature & 1.0 \\
-- & LoRA rank / $\alpha$ & 16 / 32 \\
-- & learning rate & $10^{-6}$ \\
\bottomrule
\end{tabular}
\caption{$^{\dagger}$ Provisional. $d$ should be calibrated by regressing the
probability that a viewpoint cluster is satisfied on the number of words the
response devotes to it, using the benchmark's own human-rated reference
responses --- a fit that requires no judge model.}
\end{table}

\section{Group-diversity objective (proposed; gated)}\label{sec:group}

The largest measured effect in this project is across samples: union coverage
over separately generated answers is 0.66--0.69, against about 0.51 for any single
answer. The objective below rewards each rollout for viewpoints the rest of its
group did not express. It uses no targets and no judge, so it applies to any
training question, and it replaces \eqref{eq:reward} in \eqref{eq:adv} unchanged.

\paragraph{Pool.} For a group $\{y_i\}_{i=1}^G$, all units of all rollouts are
greedily clustered by cosine similarity at threshold $\tau_p$ into a pool
$\mathcal{K}$. A unit joins its best cluster if the match clears $\tau_p$,
otherwise it opens a new one. No minimum support is imposed: a viewpoint voiced
by a single rollout is kept.

\paragraph{Covered set.} Rollout $i$ articulates cluster $k$ if some unit matches
it and at least $d_p$ words of its units are attributed to $k$, by the same
mention and depth rules as \eqref{eq:covered}. Write $C_i\subseteq\mathcal{K}$.

\paragraph{Reward.}
\begin{equation}\label{eq:group}
  q_i=\frac{|C_i|}{|\mathcal{K}|},\qquad
  n_i=\frac{\big|C_i\setminus\bigcup_{j\neq i}C_j\big|}{|\mathcal{K}|},\qquad
  r_i=q_i\,(1+\lambda\,n_i).
\end{equation}
$n_i$ is the rollout's leave-one-out contribution to the group's union; $\lambda=0$
recovers within-answer coverage of the pool. In a \emph{reference} variant, $n_i$
instead removes clusters covered by samples of the frozen base model for the same
prompt, which credits a viewpoint that several rollouts discover together.
Neither variant can separate identical rollouts, so a fully collapsed group still
gives zero advantage.

\paragraph{Gate.} Training is conditional on a within-question test against the
judge. For rollouts $i$ of one question, with $C^J_i$ the clusters the judge marks
as covered, the judge target is
$t_i=\big|C^J_i\setminus\bigcup_{k\neq i}C^J_k\big|+|C^J_i|$. Let $\kappa$ be the
pairwise concordance of $r$ with $t$, and $\kappa_J$ the concordance of the judge
with itself when re-run on a different participant subset. Variants
$(\lambda,d_p,\text{mode})$ are chosen on a random half of the questions; on the
other half the reward passes iff $\kappa_J>0.5$,
$\kappa\ge 0.5+0.5(\kappa_J-0.5)$, the bootstrap lower bound of $\kappa$ exceeds
$0.5$, and concordance with $|C^J_i|$ alone is at least $0.5$.

\paragraph{To log per group (not yet implemented).} Reward standard deviation, the
fraction of zero-variance groups, $|\mathcal{K}|$, and the mean novelty $\bar n$,
which separate a group that has collapsed from one whose reward carries no signal.

\section*{Open issues}

\begin{enumerate}
\item \textbf{Reward--metric agreement has been measured, and fails.} The
advantage \eqref{eq:adv} depends only on the \emph{within-group ordering} of
rewards, so the relevant statistic is within-question pairwise concordance with
the evaluation judge. It is $0.152$ at the defaults ($\tau = 0.50$, $d = 60$)
and at best $0.301$ across a sweep of $\tau \in [0.25, 0.50]$ and
$d \in [0, 150]$, against chance $0.500$. The failure has two regimes and no
setting escapes both: at low $\tau$ few pairs tie but the reward orders them
wrongly (concordance among separated pairs $0.417$), while at high $\tau$ $73\%$
of pairs tie. Recalibrating $\tau$ or $d$ therefore cannot repair it, and
training on \eqref{eq:reward} as specified is not justified. The judge itself is
not validated at this level (aggregate $\rho = 0.167$ against the benchmark
authors' $0.88$, ceiling unmeasured), so the concordance target is noisy too.
\item \textbf{Reachability.} \eqref{eq:grpo} reweights the reference policy's
existing distribution; it cannot create behaviour the base model never samples.
A best-of-$K$ sweep on the frozen base establishes whether the target behaviour
is reachable at all.
\item \textbf{Construct gap --- the leading explanation for item 1.} Positions
are survey answer options aggregated over demographic subgroups, about $19.9$ per
question. The evaluation's viewpoint clusters, about $7.75$ per question, are
formed by $k$-means over participants' agree/disagree votes on one another's
written statements. The two differ in both kind and granularity, so a response
can cover many positions while covering few clusters. Re-targeting the reward at
cluster-shaped objectives is the direct test and has not yet completed.
\item \textbf{Calibration of $d$.} See table footnote; the reward's absolute
scale is not meaningful until $d$ is fit.
\end{enumerate}

\end{document}
